\documentclass{article}
\usepackage[T2A]{fontenc}
\usepackage[utf8]{inputenc}
\usepackage[russian,english]{babel}
\usepackage{graphicx}
\usepackage{float}
\usepackage[hidelinks]{hyperref}
\usepackage{fancyvrb}
\usepackage{array}
\usepackage{booktabs}
\usepackage{siunitx}
\usepackage{tabularx}
\usepackage{longtable}
\usepackage{threeparttable}
\usepackage{natbib}
\usepackage[style=numeric]{biblatex}
\addbibresource{learnlatex.bib} file of reference info



\usepackage[T1]{fontenc}
\begin{document}
Hey world!
This is a first document.
\end{document}




\title{Computer Skills for Scientific Writing\\Работа с библиографией в LaTeX}
\author{Абрамян Артём Арменович}
\date{Октябрь 2026}

\begin{document}
  \maketitle
  \pagebreak

  \textbf{Цель работы:} Изучить структуру документа в LaTeX. \smallskip

  \section{Введение}
  \label{sec:intro}


  В данной работе мы изучаем структура документа в LaTeX согласно документации.


  \section{Базовая структура}
  \label{sec:basic}

  Базовая структура документа LaTeX

  Hey world!
  This is a first document.

  \smallskip

  \section{Комментарии}
  \label{sec:сomments}

  %This is a comment
  This is a simple document\footnote{with a footnote}.

  This is a new paragraph.

  \section{Запуск LaTeX}
  \label{sec:running}

  Мы можем преобразовать наш латех документ в pdf используя комманду pdflatex lab2.tex

  \section{Специальные символы}
  \label{sec:chars}

  Если вам нужно ввести специальный символ, в большинстве случаев вы можете просто поставить перед ним обратную косую черту, например,
   \{

  \section{Упражнения}
  \label{sec:ex}

  \subsection{Попробуйте добавить текст в свой первый документ, отформатировать его и посмотреть, как он изменится в вашем PDF-файле. Создайте несколько разных абзацев и добавьте переменные пробелы. Изучите, как работает ваш редактор; щелкните по исходному коду и найдите, как перейти к той же строке в вашем PDF-файле.
Попробуйте добавить несколько жестких пробелов и посмотрите, как они влияют на перенос строк.}
  \label{subsec:first}

  This is a simple document\footnote{with a footnote}.

  This is a new paragraph.
  This is hard~space

\end{document}


  \section{Перенос информации из базы данных}
  \label{sec:overflow}

  Чтобы добавить информацию в документ, нужно выполнить три шага. Во-первых, используйте LaTeX для компиляции документа, что создаст файл со списком ссылок,
  на которые ссылается ваш документ. Во-вторых, запустите программу, которая извлекает информацию из базы данных,
  выбирает используемые вами ссылки и упорядочивает их. Наконец,
  снова скомпилируйте документ, чтобы LaTeX мог использовать эту информацию для
  обработки ваших ссылок. Обычно требуется как минимум две компиляции для обработки всех
  ссылок.
  Для второго шага широко используются две системы: BibTeX и Biber. Biber
  используется только с пакетом LaTeX под названием biblatex, тогда как BibTeX используется
  либо без пакетов, либо с natbib.
  Запуск второго инструмента, а также LaTeX, обрабатывается по-разному в разных
  редакторах. Для наших онлайн-примеров есть несколько «закулисных» скриптов, которые
  делают всё за один раз. В вашем редакторе может быть всего одна кнопка «Выполнить», или вам
  возможно, придётся вручную запускать BibTeX или Biber между запусками LaTeX.
  Формат ссылок и библиографических ссылок не зависит от вашей базы данных BibTeX и
  задаётся так называемым «стилем». Мы увидим, что они работают немного по-разному в
  рабочем процессе BibTeX и biblatex, но общая идея остаётся неизменной: мы можем выбирать,
  как будут выглядеть ссылки.



  \section{Рабочий процесс BibTeX с natbib}
  \label{sec:overflowfix}


  Хотя вставлять ссылки в документ LaTeX можно и без загрузки каких-либо пакетов, возможности этого довольно ограничены. Вместо этого мы воспользуемся пакетом natbib, который
  позволяет создавать различные типы ссылок и имеет множество доступных стилей.
  Базовая структура наших входных данных показана в этом примере.

  \smallskip

  \begin{center}
    The mathematics showcase is from \citet{Graham1995}, whereas
    there is some chemistry in \citet{Thomas2008}.
    Some parenthetical citations: \citep{Graham1995}
    and then \citep[p.~56]{Thomas2008}.
    \citep[See][pp.~45--48]{Graham1995}
    Together \citep{Graham1995,Thomas2008}
  \bibliographystyle{plainnat}
  \bibliography{learnlatex}
  \end{center}




  \section{Рабочий процесс biblatex}
  \label{sec:short}

  Пакет biblatex работает несколько иначе, чем natbib, поскольку мы выбираем
базы данных в преамбуле, но печатаем их в теле документа. Для этого есть несколько
новых команд.

  \smallskip

\begin{center}
The mathematics showcase is from \autocite{Graham1995}.
Some more complex citations: \parencite{Graham1995} or
\textcite{Thomas2008} or possibly \citetitle{Graham1995}.
\autocite[56]{Thomas2008}
\autocite[See][45-48]{Graham1995}
Together \autocite{Thomas2008,Graham1995}
\printbibliography
\end{center}

  \smallskip

  Обратите внимание, что addbibresource требует полного имени файла базы данных, тогда как мы
опустили расширение .bib для bibliography с natbib. Также обратите внимание, что biblatex
использует довольно длинные имена для своих команд цитирования, но все они довольно легко
угадываются.
Опять же, короткий текст до и после ссылки можно вставить с помощью необязательных
аргументов. Обратите внимание, что номера страниц не обязательно должны иметь префикс p.~ или pp.~,
biblatex может автоматически добавить соответствующий префикс.
В biblatex стиль ссылки выбирается при загрузке пакета. Здесь мы
использовали authoryear, но есть и числовой стиль, а также множество других.



  \section{Выбор между BibTeX и BibLaTeX}
  \label{sec:rules}


  Несмотря на то, что и рабочий процесс BibTeX, и biblatex получают входные данные через BibTeX-файлы и могут создавать структурно схожие выходные данные в документе, они используют совершенно разные способы получения этого результата. Это означает, что между двумя подходами есть некоторые различия, которые могут помочь вам выбрать наиболее подходящий для вас.
В рабочем процессе BibTeX стиль библиографии в конечном итоге определяется файлом .bst,
который вы выбираете с помощью команды bibliographystyle. biblatex не
использует файлы .bst и использует другую систему. Если вы используете шаблон, который поставляется с файлом .bst или вам предоставлен файл .bst для вашего проекта, вы должны использовать рабочий процесс BibTeX и не можете использовать biblatex.
Различный подход, используемый biblatex, подразумевает, что вы можете изменять вывод команд библиографии и цитирования непосредственно из преамбулы документа с помощью команд на основе LaTeX. С другой стороны, модификации стилей BibTeX .bst
обычно требуют работы с этими внешними файлами и знания
языка программирования BibTeX. Как правило, biblatex считается
проще в настройке, чем рабочий процесс BibTeX.
В biblatex, как правило, проще реализовать более сложные стили цитирования с
более широким набором различных команд цитирования. Он также предлагает больше контекстно-зависимых
функций. Грубо говоря, это менее интересно для стилей, распространённых во многих
предметах STEM, но становится актуальным для некоторых более сложных стилей в некоторых областях
гуманитарных наук.
BibTeX может корректно сортировать только символы US-ASCII и использует обходные пути для
обеспечения сортировки на основе US-ASCII для символов, не входящих в US-ASCII. Благодаря Biber, biblatex предлагает все возможности сортировки в Unicode. Таким образом, biblatex обычно является лучшим
выбором, если вы хотите отсортировать библиографию в порядке «не ASCII/не английский». Рабочий процесс BibTeX, существующий гораздо дольше, чем biblatex, более
устоявшийся, чем biblatex, а это означает, что многие издательства и журналы ожидают
библиографии, созданные с помощью рабочего процесса BibTeX. Эти издательства не могут или, как правило, не принимают материалы, созданные с помощью biblatex.
Суть в следующем: проверьте правила для авторов/подачи материалов, если вы отправляете
работу в журнал или издательство. Если вам предоставлен файл .bst, вам необходимо использовать рабочий процесс BibTeX. Если вам нужен относительно простой стиль библиографии и цитирования и
нужна сортировка только на основе английской кодировки US-ASCII, рабочего процесса BibTeX должно быть достаточно. Если вам
нужен более сложный стиль цитирования, сортировка не на английском языке или более простой доступ к
функциям настройки стиля цитирования и библиографии, вам стоит рассмотреть возможность использования biblatex.


  \section{Работа с сортировкой, отличной от английского}
  \label{sec:merge}

  Программа BibTeX была написана в первую очередь для работы со ссылками на английском языке. Она
очень ограничена в обработке диакритических знаков и ещё более ограничена в работе с нелатинскими буквами. В отличие от этого, программа Biber изначально была написана для корректной работы
со смешанным алфавитом.
Это означает, что если вы сортируете библиографию и вам нужно отсортировать её в порядке, отличном от английского, вам действительно следует использовать biblatex и Biber,
а не natbib и BibTeX.








  \section{Гиперссылки}
  \label{sec:other}


  Если вы загрузите пакет hyperref, он автоматически преобразует часть контента вашей библиографии в ссылки. Это особенно полезно для URL-адресов и DOI.






  \section{Различия в практике ввода BibTeX}
  \label{sec:ell}

  Хотя общий синтаксис файлов BibTeX одинаков, независимо от того, используете ли вы рабочий процесс BibTeX или biblatex, набор поддерживаемых полей (используемых стилем) и их точное значение могут различаться не только между рабочим процессом BibTeX и biblatex,
но и между разными стилями BibTeX. Большой «основной набор» типов записей и полей
одинаков почти для всех стилей, но в некоторых полях есть различия.
Распространённым примером является URL. Некоторые старые стили BibTeX .bst (в первую очередь «стандартные стили BibTeX», например, plain.bst, unsrt.bst и т. д.) появились ещё до изобретения URL и не имеют отдельного поля для URL-адреса онлайн-ресурса. Многие более новые
стили имеют отдельное поле URL. Обходной путь для отображения URL в старых
стилях обычно заключается в использовании поля howpublished, но в новых стилях, конечно,
предпочтительнее использовать специальное поле URL.
Чтобы использовать весь потенциал используемого вами стиля, вам нужно
узнать, какой набор полей он поддерживает и каково их семантическое значение.

  \smallskip














  \section{Упражнения}
  \label{sec:ex}

  \subsection{Add a citation that’s not in the database and see how it appears.}
  \label{subsec:width}

\begin{center}
The mathematics showcase is from \autocite{Graham1995}.
Some more complex citations: \parencite{Graham1995} or
\textcite{checkMissingCite} or possibly \citetitle{Graham1995}.
\autocite[56]{Thomas2008}
\autocite[See][45-48]{Graham1995}
Together \autocite{Thomas2008,Graham1995}
\printbibliography
\end{center}



  \subsection{Experiment with natbib’s numeric and biblatex’s style=numeric option.}
  \label{subsec:width}


  \section{Заключение}
  \label{sec:conclusion}

  В данной работе были продемонстрированы:
  \begin{itemize}
    \item Базовое включение таблицы в документ (раздел~\ref{sec:basic})
    \item Переполнение таблицы (раздел~\ref{sec:overflow})
    \item Фикс проблемы переполнение таблицы (раздел~\ref{sec:overflowfix})
    \item Сокращенное задание типов колон (раздел~\ref{sec:short})
    \item Добавление правил (раздел~\ref{sec:rules})
    \item Слияние столбцов (раздел~\ref{sec:merge})
    \item Другое (раздел~\ref{sec:other})
    \item Настройка правил BookTabs (раздел~\ref{sec:booktabsrules})
    \item Числовое выравнивание в колоннах (раздел~\ref{sec:alignment})
    \item Общая ширина таблицы (раздел~\ref{sec:tablewidth})
    \item Определение новых типов колон (раздел~\ref{sec:newtypes})
    \item Vertikal tricks (раздел~\ref{sec:vertikaltricks})
  \end{itemize}

  \textbf{Выводы:} В данной работе мы изучили возможности работы включения и манипулирования таблицами в LaTeX, освоили различные типы колон, научились использовать правила.

\end{document}
